% !PRINTED 647
\weberpage{647}
\begin{center}
\S.\,187.

\gesperrt{Realit\"atsverh\"altnisse.}
\end{center}

Wenn der Körper $\mathfrak{K}$ ein reeller ist, so giebt es, wie wir im §.\,180 gesehen haben, zwei Arten von metacyklischen Gleichungen von Primzahlgrad $n$, solche mit einer reellen und $n-1$ imaginären Wurzeln und solche mit lauter reellen Wurzeln. Ebenso haben wir im §.\,165 gesehen, dass es zwei Arten von cyklischen Gleichungen eines geraden Grades giebt, nämlich solche mit lauter reellen und solche mit lauter imaginären Wurzeln.

% [sic] Print runs both the $k$ and $\tau$ lists through index $n-1$ here;
% the cyclic equation and radical system defined in §.185 run through $n-2$.
Wenn nun die cyklische Gleichung, deren Wurzeln die Grössen $k_0, k_1 \ldots k_{n-1}$ sind, zur ersten Art gehört, wenn also $k_0, k_1 \ldots k_{n-1}$ reell sind und die Radicale $\tau_0, \tau_1, \ldots \tau_{n-1}$ auch reell genommen werden, so zeigt die Formel (10) §.\,186, dass $\xi_0$ reell, $\xi_h$ und $\xi_{-h}$ conjugirt imaginär sind.

Wenn andererseits die $k$ imaginär sind, so ist nach §.\,165 $k_\nu$ und $k_{\nu+\frac{n-1}{2}}$ conjugirt imaginär, und nach der Formel (8) in §.\,185 können wir setzen, wenn $\Phi$ eine rationale Function bedeutet,

\[
\sqrt[n]{R_{\nu+\frac{n-1}{2}}} = \Phi(k_\nu) \left(\sqrt[n]{R_\nu}\right)^{g^{\frac{n-1}{2}}},
\]
oder auch

\[
\sqrt[n]{R_{\nu+\frac{n-1}{2}}}\ \sqrt[n]{R_\nu} = \Phi(k_\nu) \left(\sqrt[n]{R_\nu}\right)^{g^{\frac{n-1}{2}}+1}.
\]

Weil nun $g$ eine primitive Wurzel von $n$ ist, so ist $g^{\frac{n-1}{2}}+1$ durch $n$ theilbar, und die rechte Seite wird daher rational. Bedeutet also $\Psi$ eine rationale Function, so ist

\[
\sqrt[n]{R_{\nu+\frac{n-1}{2}}}\ \sqrt[n]{R_\nu} = \Psi(k_\nu),
\]
und diese Formel zeigt, dass

\[
\Psi(k_\nu) = \Psi\!\left(k_{\nu+\frac{n-1}{2}}\right),
\]
also dass $\Psi(k_\nu)$ reell ist. Daraus folgt, dass

\[
\sqrt[n]{R_{\nu+\frac{n-1}{2}}} \quad \text{und} \quad \sqrt[n]{R_\nu}
\]
%CONT>>
